\section{Week 6}
\subsection{D.3.1 Finite Domain (1D) (Series Expansion)}

\subsubsection{D.3.1.1 Fixed Concentration Boundaries}

The governing Partial Differential Equation (PDE) for 1D time-dependent diffusion is:
\begin{equation}
\frac{\partial C}{\partial t} = D \frac{\partial^{2}C}{\partial x^{2}}
\end{equation}

\paragraph{Boundary Conditions (BC)}
We assume fixed, zero concentration at the physical boundaries $x=0$ and $x=L$:
\begin{equation}
C(0,t) = C(L,t) = 0
\end{equation}

\begin{conceptbox}[Superposition Principle]
The PDE is \textbf{linear} and \textbf{homogeneous}, and it is subject to \textbf{Homogeneous Boundary Conditions (HBC)}. Therefore, if $C_{1}(x,t)$ and $C_{2}(x,t)$ are valid solutions, any linear combination is also a valid solution:
\[
C_{1}(x,t) + C_{2}(x,t) \quad \text{and} \quad \beta C_{1}(x,t) \quad (\text{for } \beta \in \mathbb{R})
\]
\textbf{Proof via PDE Linearity:}
\begin{align*}
\frac{\partial}{\partial t}(C_{1}+C_{2}) - D\frac{\partial^{2}}{\partial x^{2}}(C_{1}+C_{2}) &= \left(\frac{\partial C_{1}}{\partial t} - D\frac{\partial^{2}C_{1}}{\partial x^{2}}\right) + \left(\frac{\partial C_{2}}{\partial t} - D\frac{\partial^{2}C_{2}}{\partial x^{2}}\right) \\
&= 0 + 0 = 0
\end{align*}
\textbf{Proof via BCs:}
\[
(C_{1}+C_{2})(0,t) = C_{1}(0,t) + C_{2}(0,t) = 0 + 0 = 0
\]
\end{conceptbox}

---

\subsubsection{Solution Approach (Method of Separation of Variables)}

Because of linearity, the general solution is constructed as an infinite linear combination of base eigenfunctions $\varphi_{n}(x,t)$:
\begin{equation}
C(x,t) = \sum_{n=1}^{\infty} A_{n} \varphi_{n}(x,t)
\end{equation}

We assume the fundamental solution $\varphi(x,t)$ can be factored into purely spatial and purely temporal functions:
\begin{equation}
\varphi(x,t) = X(x)T(t)
\end{equation}

Substituting this ansatz into the PDE gives $X(x)T'(t) = D X''(x)T(t)$. Dividing by $D X(x)T(t)$ separates the variables:
\[
\frac{T'(t)}{D T(t)} = \frac{X''(x)}{X(x)} = -\lambda
\]
Since the LHS strictly depends on $t$ and the RHS strictly depends on $x$, both must equal a separation constant, denoted $-\lambda$. This yields two coupled Ordinary Differential Equations (ODEs):
\begin{align}
T'(t) + D\lambda T(t) &= 0 \quad \text{(Time dependence)} \\
X''(x) + \lambda X(x) &= 0 \quad \text{(Spatial dependence)}
\end{align}

---

\paragraph{Solving the Spatial Eigenvalue Problem}
We must solve the $2^{\text{nd}}$ order ODE $X''(x) + \lambda X(x) = 0$ subject to the HBCs $X(0) = X(L) = 0$.

\begin{itemize}
    \item \textbf{Case I ($\lambda < 0$):} $X(x) = \alpha_{1}e^{\sqrt{-\lambda}x} + \alpha_{2}e^{-\sqrt{-\lambda}x}$. Applying BCs gives $\alpha_1 = \alpha_2 = 0$. Trivial solution.
    \item \textbf{Case II ($\lambda = 0$):} $X(x) = \alpha_{1}x + \alpha_{2}$. Applying BCs gives $\alpha_2=0$ and $\alpha_1 L = 0$. Trivial solution.
    \item \textbf{Case III ($\lambda > 0$):} The general solution is a harmonic oscillator form:
    \[
    X(x) = \alpha_{1}\cos(\sqrt{\lambda}x) + \alpha_{2}\sin(\sqrt{\lambda}x)
    \]
    Applying $X(0) = 0 \implies \alpha_{1} = 0$. \\
    Applying $X(L) = 0 \implies \alpha_{2}\sin(\sqrt{\lambda}L) = 0$. \\
    To avoid the trivial solution ($\alpha_{2} \neq 0$), the sine term must vanish: $\sin(\sqrt{\lambda}L) = 0$.
\end{itemize}

This yields the \textbf{eigenvalues} and spatial \textbf{eigenfunctions}:
\begin{eqbox}{Spatial Eigenfunctions}
\begin{align*}
\sqrt{\lambda}L = n\pi \quad &\implies \quad \lambda_{n} = \left(\frac{n\pi}{L}\right)^{2} \quad \text{for } n \in \mathbb{N}^* \\
X_{n}(x) &= \sin\left(\frac{n\pi}{L}x\right)
\end{align*}
\textit{Note: The constant $\alpha_{n}$ is omitted as it will be absorbed by the Fourier coefficients $A_{n}$.}
\end{eqbox}

---

\paragraph{Solving the Temporal Problem}
For a given eigenvalue $\lambda_n$, the temporal ODE is a first-order linear decay:
\[
T_{n}'(t) = -D\lambda_{n}T_{n}(t) \implies T_{n}(t) = \beta_{n}e^{-D\lambda_{n}t}
\]
Substituting $\lambda_n$, the temporal amplitude decays exponentially:
\begin{equation}
T_{n}(t) = e^{-D\left(\frac{n\pi}{L}\right)^{2}t}
\end{equation}

---

\subsubsection{Combining the Solutions \& Initial Conditions}

Multiplying the spatial and temporal components, the general solution satisfying the boundary conditions is:
\begin{equation}
C(x,t) = \sum_{n=1}^{\infty} A_{n} e^{-D\left(\frac{n\pi}{L}\right)^{2}t} \sin\left(\frac{n\pi}{L}x\right)
\end{equation}

\paragraph{Determining the Fourier Coefficients $A_n$}
At $t=0$, the solution must match the given initial concentration profile $C_{0}(x)$:
\[
C(x,0) = C_{0}(x) = \sum_{n=1}^{\infty} A_{n}\sin\left(\frac{n\pi}{L}x\right)
\]
To isolate $A_n$, we multiply both sides by $\sin\left(\frac{m\pi}{L}x\right)$ (for integer $m$) and integrate over the domain $[0, L]$:
\[
\int_{0}^{L} C_{0}(x)\sin\left(\frac{m\pi}{L}x\right) dx = \sum_{n=1}^{\infty} A_{n} \int_{0}^{L} \sin\left(\frac{n\pi}{L}x\right)\sin\left(\frac{m\pi}{L}x\right) dx
\]

\begin{expertbox}[Fourier Orthogonality]
Using the trigonometric identity $\sin(a)\sin(b) = \frac{1}{2}[\cos(a-b) - \cos(a+b)]$, the integral of the product of two sine waves over the domain evaluates to:
\[
\int_{0}^{L} \sin\left(\frac{n\pi}{L}x\right)\sin\left(\frac{m\pi}{L}x\right) dx = 
\begin{cases} 
0 & \text{if } n \neq m \\ 
\frac{L}{2} & \text{if } n = m 
\end{cases}
\]
\end{expertbox}

Due to this orthogonality, the infinite sum collapses to a single non-zero term when $n=m$:
\[
\int_{0}^{L} C_{0}(x)\sin\left(\frac{n\pi}{L}x\right) dx = A_{n}\left(\frac{L}{2}\right) \implies A_{n} = \frac{2}{L}\int_{0}^{L} C_{0}(x)\sin\left(\frac{n\pi}{L}x\right) dx
\]

\begin{resultboxnamed}{Complete Analytical Solution}
For an arbitrary initial condition $C_0(x)$, the full space-time solution is the Fourier sine series:
\begin{equation}
C(x,t) = \sum_{n=1}^{\infty} \underbrace{\left[ \frac{2}{L}\int_{0}^{L} C_{0}(\xi)\sin\left(\frac{n\pi}{L}\xi\right) d\xi \right]}_{\text{Fourier Coefficient } A_n} \cdot \underbrace{e^{-D\left(\frac{n\pi}{L}\right)^{2}t}}_{\text{Temporal Decay}} \cdot \underbrace{\sin\left(\frac{n\pi}{L}x\right)}_{\text{Spatial Mode}}
\end{equation}
\textit{Note: The dummy integration variable has been set to $\xi$ to strictly separate it from the spatial coordinate $x$.}
\end{resultboxnamed}